\documentclass[a4paper,12pt]{article}
\usepackage{fontspec} % Пакет для работы со шрифтами в XeLaTeX
\setmainfont{Liberation Serif} % Использование стандартного шрифта
\usepackage[russian]{babel}

% Математические пакеты
\usepackage{amsmath, amsfonts, amssymb}
\usepackage{geometry}
\usepackage{graphicx}
\usepackage{tabularx}
\usepackage{booktabs}
\usepackage{indentfirst}
\usepackage{pgfplots}
\pgfplotsset{compat=1.18}

\geometry{top=2cm, bottom=2cm, left=2.5cm, right=1.5cm}

\begin{document}

\thispagestyle{empty}

% Шапка: Название университета и факультет
\begin{center}
    {\fontsize{14}{16}\selectfont \textbf{Университет ИТМО}} \\
    {\fontsize{12}{14}\selectfont Физико-технический мегафакультет} \\
    {\fontsize{12}{14}\selectfont Физический факультет}
\end{center}

% Место для логотипа (справа сверху)
\begin{flushright}
    \vspace{-3cm}
    \Huge \textbf{ИТМО} % Плейсхолдер логотипа
\end{flushright}

\vspace{1.5cm}

% Блок с информацией о студентах и проверке
\noindent
\begin{tabular}{ll}
    Поток \makebox[5cm]{ФИЗ-3 Э БИТ 1.2.1} & К работе допущен\underline{\makebox[3cm]{}} \\[0.3cm]
    Студенты: {\makebox[5.5cm]{Пахомов Владимир Юрьевич,}} & Работа выполнена\underline{\makebox[3cm]{}} \\
    Сивков Иван Александрович & \\
    Преподаватель: Терещенко Георгий Викторович & Отчет принят\underline{\makebox[3.5cm]{}}
\end{tabular}

\vspace{2cm}

% Заголовок работы
\begin{center}
    {\fontsize{18}{22}\selectfont \textbf{Рабочий протокол и отчет}} \\[0.5cm]
    {\fontsize{18}{22}\selectfont \textbf{по лабораторной работе № 5.07}} \\[1.5cm]
    
    {\fontsize{16}{20}\selectfont \textbf{Определение постоянной Планка методом}} \\
    {\fontsize{16}{20}\selectfont \textbf{задерживающего потенциала}}
\end{center}

\newpage

\section{Цель работы}
1) Экспериментально проверить законы фотоэффекта. \\
2) Определение постоянной Планка и работы выхода электрона из металла.

\section{Задачи, решаемые при выполнении работы}
1) Определение запирающего напряжения; \\
2) Изучение зависимости кинетической энергии электронов от частоты падающего излучения; \\
3) Экспериментальное определение материала фотокатода и вычисление постоянной Планка.

\section{Объект исследования}
Вакуумный фотоэлемент на основе цезиевого фотокатода.

\section{Метод исследования}
Метод основан на измерении задерживающего напряжения $U_0$ (противо-ЭДС), необходимого для прекращения фототока при освещении катода монохроматическим светом различных частот. На основе уравнения Эйнштейна для фотоэффекта строится линейная зависимость, угловой коэффициент которой позволяет рассчитать постоянную Планка $h$, а свободный член — работу выхода электронов $A_{\text{вых}}$.

\section{Рабочие формулы и исходные данные}

\begin{equation}
h\nu = A_{\text{вых}} + T \tag{1}
\end{equation}
— уравнение Эйнштейна для внешнего фотоэффекта.

\begin{equation}
T = eU_0 \tag{2}
\end{equation}
— кинетическая энергия фотоэлектрона через запирающее напряжение $U_0$.

\begin{equation}
U_0 = \frac{h}{e}\nu - \frac{A_{\text{вых}}}{e} \tag{3}
\end{equation}
— зависимость запирающего напряжения от частоты падающего излучения.

\begin{equation}
\nu = \frac{c}{\lambda} \tag{4}
\end{equation}
— связь частоты $\nu$ с длиной волны $\lambda$ ($c = 2{,}99 \cdot 10^8$ м/с).

\begin{equation}
\nu_0 = \frac{A_{\text{вых}}}{h} \tag{5}
\end{equation}
— частота красной границы фотоэффекта.

\vspace{0.5cm}
\textbf{Константы:}
\begin{itemize}
    \item Заряд электрона: $e = 1{,}602 \cdot 10^{-19}$ Кл
    \item Скорость света в вакууме: $c = 2{,}99 \cdot 10^8$ м/с
\end{itemize}

\newpage

\section{Результаты прямых измерений}

Результаты измерений длин волн светодиодов и соответствующих задерживающих напряжений, а также рассчитанные на их основе частоты и энергии, занесены в таблицу 1.

\begin{center}
\begin{tabular}{|c|c|c|c|}
\hline
$\lambda$, нм & $U_0$, В & $\nu$, ТГц & $T = eU_0$, $\cdot 10^{-19}$ Дж \\ \hline
472 & 0,633 & 633,47 & 1,0141 \\ \hline
505 & 0,486 & 592,08 & 0,7786 \\ \hline
525 & 0,437 & 569,52 & 0,7001 \\ \hline
588 & 0,143 & 508,50 & 0,2291 \\ \hline
611 & 0,081 & 489,36 & 0,1298 \\ \hline
\end{tabular}

\medskip
Таблица 1 — Результаты измерений и первичных вычислений.
\end{center}

\section{Обработка результатов}

Выполним обработку полученных данных строго по пунктам методических указаний:

\subsection*{Пункт 1. Построение графика зависимости энергии электронов $T(\nu)$}

Ниже представлен график зависимости максимальной кинетической энергии фотоэлектронов $T$ от частоты падающего излучения $\nu$.

\includegraphics[width=17cm]{graphic.png}

\subsection*{Пункт 2. Линейная аппроксимация методом наименьших квадратов}

Для аппроксимации зависимости вида $T = k\nu + b$ используем метод наименьших квадратов (МНК). Частоту $\nu$ подставим в Гц, энергию $T$ — в Дж.
Вычислим промежуточные суммы ($n = 5$):
\begin{itemize}
    \item $\sum \nu_i = 2792{,}93 \cdot 10^{12}$ Гц
    \item $\sum T_i = 2{,}8517 \cdot 10^{-19}$ Дж
    \item $\sum \nu_i^2 = 1{,}57424 \cdot 10^{30}$ Гц$^2$
    \item $\sum (\nu_i T_i) = 1{,}68213 \cdot 10^{-5}$ Дж$\cdot$Гц
\end{itemize}

Угловой коэффициент (коэффициент наклона) аппроксимирующей прямой:
\begin{equation*}
k = \frac{n \sum (\nu_i T_i) - \sum \nu_i \sum T_i}{n \sum \nu_i^2 - (\sum \nu_i)^2} \approx 6{,}305 \cdot 10^{-34} \text{ Дж}\cdot\text{с}
\end{equation*}

Свободный член аппроксимирующей прямой:
\begin{equation*}
b = \frac{\sum T_i - k \sum \nu_i}{n} \approx -2{,}952 \cdot 10^{-19} \text{ Дж}
\end{equation*}

\subsection*{Пункт 3. Определение угла наклона, красной границы и их погрешностей}

1. **Угол наклона (тангенс угла наклона) аппроксимирующей прямой:**
Согласно теории (стр. 5 методического пособия), тангенс угла наклона прямой зависимости $T(\nu)$ численно равен постоянной Планка $h$:
\begin{equation*}
\text{tg}(\alpha) = k = h = 6{,}305 \cdot 10^{-34} \text{ Дж}\cdot\text{с}
\end{equation*}
Погрешность коэффициента наклона (постоянной Планка) рассчитана ниже в Пункте 5:
\begin{equation*}
\Delta (\text{tg}\,\alpha) = \Delta h = 0{,}032 \cdot 10^{-34} \text{ Дж}\cdot\text{с}
\end{equation*}

2. **Частота красной границы фотоэффекта ($\nu_0$):**
Из условия прекращения фотоэффекта ($T = 0 \implies h\nu_0 = A_{\text{вых}} = -b$):
\begin{equation*}
\nu_0 = \frac{A_{\text{вых}}}{h} = \frac{2{,}952 \cdot 10^{-19} \text{ Дж}}{6{,}305 \cdot 10^{-34} \text{ Дж}\cdot\text{с}} \approx 4{,}682 \cdot 10^{14} \text{ Гц} = 468{,}2 \text{ ТГц}
\end{equation*}

Определим погрешность частоты красной границы $\Delta \nu_0$:
\begin{equation*}
\frac{\Delta \nu_0}{\nu_0} = \sqrt{\left(\frac{\Delta A}{A}\right)^2 + \left(\frac{\Delta h}{h}\right)^2} = \sqrt{0{,}0052^2 + 0{,}0051^2} \approx 0{,}0073 \quad (0{,}73\%)
\end{equation*}
\begin{equation*}
\Delta \nu_0 = \nu_0 \cdot 0{,}0073 = 468{,}2 \text{ ТГц} \cdot 0{,}0073 \approx 3{,}4 \text{ ТГц}
\end{equation*}

Таким образом, частота красной границы составляет:
\begin{equation*}
\nu_0 = (468 \pm 3) \text{ ТГц}
\end{equation*}

\subsection*{Пункт 4. Определение материала фотокатода}

Экспериментальное значение работы выхода электронов из фотокатода:
\begin{equation*}
A_{\text{вых}} = -b = 2{,}952 \cdot 10^{-19} \text{ Дж} \approx 1{,}843 \text{ эВ}
\end{equation*}

Согласно таблицам физических величин, полученная работа выхода в пределах погрешности идеально соответствует **цезию (Cs)**, справочное значение работы выхода которого составляет $A_{\text{вых}}^{\text{теор}} = 1{,}89$ эВ. Относительное отклонение экспериментального значения от теории составляет всего $2{,}6\%$.

\subsection*{Пункт 5. Расчет погрешностей постоянной Планка и работы выхода}

В соответствии с указаниями к работе, расчет погрешностей выполняется для длины волны $\lambda = 500$ нм, при этом относительная погрешность вольтметра составляет $0{,}5\% = 0{,}005$, а абсолютная погрешность длины волны $\Delta \lambda = 0{,}5$ нм.

1. **Относительная погрешность постоянной Планка (коэффициента наклона):**
\begin{equation*}
\frac{\Delta h}{h} = \sqrt{\left(\frac{\Delta U}{U}\right)^2 + \left(\frac{\Delta \lambda}{\lambda}\right)^2} = \sqrt{0{,}005^2 + \left(\frac{0{,}5}{500}\right)^2} \approx 0{,}0051 \quad (0{,}51\%)
\end{equation*}
Абсолютная погрешность:
\begin{equation*}
\Delta h = h \cdot 0{,}0051 = 6{,}305 \cdot 10^{-34} \cdot 0{,}0051 \approx 0{,}032 \cdot 10^{-34} \text{ Дж}\cdot\text{с}
\end{equation*}

2. **Относительная погрешность работы выхода:**
\begin{equation*}
\frac{\Delta A}{A} = \sqrt{\left(\frac{\Delta h}{h}\right)^2 + \left(\frac{\Delta \lambda}{\lambda}\right)^2} = \sqrt{0{,}0051^2 + 0{,}001^2} \approx 0{,}0052 \quad (0{,}52\%)
\end{equation*}
Абсолютная погрешность:
\begin{equation*}
\Delta A_{\text{вых}} = 2{,}952 \cdot 10^{-19} \cdot 0{,}0052 \approx 0{,}015 \cdot 10^{-19} \text{ Дж} \approx 0{,}01 \text{ эВ}
\end{equation*}

Запишем итоговые доверительные интервалы для измеряемых величин:
\begin{equation*}
h = (6{,}31 \pm 0{,}03) \cdot 10^{-34} \text{ Дж}\cdot\text{с}
\end{equation*}
\begin{equation*}
A_{\text{вых}} = (2{,}95 \pm 0{,}02) \cdot 10^{-19} \text{ Дж} = (1{,}84 \pm 0{,}01) \text{ эВ}
\end{equation*}

\section{Вывод}

В ходе лабораторной работы была экспериментально подтверждена линейная зависимость задерживающего потенциала от частоты падающего монохроматического излучения, что полностью согласуется с уравнениями внешнего фотоэффекта Эйнштейна.

Построен график зависимости $T(\nu)$, угол наклона которого (физически равный постоянной Планка) был найден с помощью метода наименьших квадратов: $h = (6{,}31 \pm 0{,}03) \cdot 10^{-34}$ Дж·с. Данное значение хорошо сходится с общепринятой фундаментальной константой Планка ($6{,}626 \cdot 10^{-34}$ Дж·с, отклонение составляет $4{,}8\%$).

Определена работа выхода электронов из фотокатода: $A_{\text{вых}} = (1{,}84 \pm 0{,}01)$ эВ, на основании чего сделан вывод о том, что материалом исследуемого катода является цезий (справочное значение $1{,}89$ эВ, отклонение составило $2{,}6\%$). Красная граница фотоэффекта для данного материала определена на частоте $\nu_0 = (468 \pm 3)$ ТГц, что соответствует длине волны $\lambda_0 \approx 639$ нм.

\vfill
\begin{flushright}
    4
\end{flushright}

\end{document}